% GENERATED FILE. Edit canonical lessons, then run npm run generate:books.
% canonical-source-hash: fnv1a64:12366df12ba7310a
% canonical-lessons: KA-C512-raktadottada, KA-C512-hrdayaghata, KA-C512-raktasrava, KA-C512-parsvavayu, KA-C512-prajne-tappu

\chapter{Serious Illness}
\label{ch:ka-serious-illness}
\hlchaptermodality{512}

\begin{chapteropening}
\textbf{By the end of this chapter:} \emph{I can say blood pressure, a heart attack, bleeding, a stroke and to faint.}

Complete the last lesson of chapter 512: I can say blood pressure, a heart attack, bleeding, a stroke and to faint.
\end{chapteropening}

\section[\texorpdfstring{\kn{ರಕ್ತದೊತ್ತಡ} (raktadottaḍa)}{ರಕ್ತದೊತ್ತಡ (raktadottaḍa)}]{\kn{ರಕ್ತದೊತ್ತಡ} (raktadottaḍa) — blood pressure}

\label{lesson:KA-C512-raktadottada}



\begin{quote}
Before the new one: say the Kannada for a symptom, then the Kannada for a vaccine.
\end{quote}

\subsection*{You'll want to know: \kn{ರಕ್ತದೊತ್ತಡ}}
\textbf{\kn{ರಕ್ತದೊತ್ತಡ}} — \emph{raktadottaḍa} — \textquotedblleft{}blood pressure\textquotedblright{}.

\begin{quote}
\kn{ಉಪ್ಪು} \kn{ಹೆಚ್ಚು} \kn{ತಿಂದರೆ} \kn{ರಕ್ತದೊತ್ತಡ} \kn{ಏರುತ್ತದೆ}.
\end{quote}

\begin{quote}
\emph{uppu heccu tindare raktadottaḍa ēruttade.}
\end{quote}

\begin{quote}
If you eat a lot of salt, your blood pressure rises.
\end{quote}

\kn{ರಕ್ತ}, blood + \kn{ಒತ್ತಡ}, pressure, joined with a \kn{ದ}.

\begin{cousinweb}
\textbf{\kn{ರಕ್ತದೊತ್ತಡ}} (\emph{raktadottaḍa}) joins Sanskrit \emph{rakta}, \textquotedblleft{}red; blood\textquotedblright{} (CDIAL 10539), and the native \emph{ottaḍa}, \textquotedblleft{}pressure\textquotedblright{}, from \emph{ottu}, \textquotedblleft{}to press\textquotedblright{} (DEDR 1021).
\end{cousinweb}

\begin{grammarlens}[title={What you've built}]
One more word for describing what you lived and felt.
\end{grammarlens}

\subsection*{Your turn}
\begin{itemize}
\raggedright
  \item \emph{Say it:} \emph{raktadottaḍa}
  \item \emph{Say it:} \emph{raktadottaḍa}, once more
  \item \emph{Say it:} \emph{uppu heccu tindare raktadottaḍa ēruttade.} — If you eat a lot of salt, your blood pressure rises.
  \item \emph{Recall:} say the Kannada for a symptom, then the Kannada for a vaccine, then say \emph{raktadottaḍa} again
\end{itemize}

\begin{tcolorbox}[breakable,colback=teal!4,colframe=teal!35!black,title={Before you move on}]
What does \kn{ರಕ್ತದೊತ್ತಡ} mean? (\textquotedblleft{}Blood pressure\textquotedblright{}.) Say it once more.
\end{tcolorbox}
\section[\texorpdfstring{\kn{ಹೃದಯಾಘಾತ} (hṛdayāghāta)}{ಹೃದಯಾಘಾತ (hṛdayāghāta)}]{\kn{ಹೃದಯಾಘಾತ} (hṛdayāghāta) — a heart attack}

\label{lesson:KA-C512-hrdayaghata}



\begin{quote}
Before the new one: say the Kannada for a vaccine, then the Kannada for blood pressure.
\end{quote}

\subsection*{You'll want to know: \kn{ಹೃದಯಾಘಾತ}}
\textbf{\kn{ಹೃದಯಾಘಾತ}} — \emph{hṛdayāghāta} — \textquotedblleft{}a heart attack\textquotedblright{}.

\begin{quote}
\kn{ಹೃದಯಾಘಾತ} \kn{ಆದ} \kn{ಮೇಲೆ} \kn{ಅಪ್ಪ} \kn{ಸಿಗರೇಟು} \kn{ಬಿಟ್ಟರು}.
\end{quote}

\begin{quote}
\emph{hṛdayāghāta āda mēle appa sigarēṭu biṭṭaru.}
\end{quote}

\begin{quote}
After his heart attack, father gave up cigarettes.
\end{quote}

\kn{ಹೃದಯ}, the heart + \kn{ಆಘಾತ}, a blow: the word for a shock, now a blow to the heart.

\begin{cousinweb}
\textbf{\kn{ಹೃದಯಾಘಾತ}} (\emph{hṛdayāghāta}) is a learned compound built from Sanskrit \emph{hṛdaya}, \textquotedblleft{}the heart\textquotedblright{}, and \emph{āghāta}, \textquotedblleft{}a blow\textquotedblright{} (Monier-Williams).
\end{cousinweb}

\begin{grammarlens}[title={What you've built}]
One more word for describing what you lived and felt.
\end{grammarlens}

\subsection*{Your turn}
\begin{itemize}
\raggedright
  \item \emph{Say it:} \emph{hṛdayāghāta}
  \item \emph{Say it:} \emph{hṛdayāghāta}, once more
  \item \emph{Say it:} \emph{hṛdayāghāta āda mēle appa sigarēṭu biṭṭaru.} — After his heart attack, father gave up cigarettes.
  \item \emph{Recall:} say the Kannada for a vaccine, then the Kannada for blood pressure, then say \emph{hṛdayāghāta} again
\end{itemize}

\begin{tcolorbox}[breakable,colback=teal!4,colframe=teal!35!black,title={Before you move on}]
What does \kn{ಹೃದಯಾಘಾತ} mean? (\textquotedblleft{}A heart attack\textquotedblright{}.) Say it once more.
\end{tcolorbox}
\section[\texorpdfstring{\kn{ರಕ್ತಸ್ರಾವ} (raktasrāva)}{ರಕ್ತಸ್ರಾವ (raktasrāva)}]{\kn{ರಕ್ತಸ್ರಾವ} (raktasrāva) — bleeding; a haemorrhage}

\label{lesson:KA-C512-raktasrava}



\begin{quote}
Before the new one: say the Kannada for blood pressure, then the Kannada for a heart attack.
\end{quote}

\subsection*{You'll want to know: \kn{ರಕ್ತಸ್ರಾವ}}
\textbf{\kn{ರಕ್ತಸ್ರಾವ}} — \emph{raktasrāva} — \textquotedblleft{}bleeding; a haemorrhage\textquotedblright{}.

\begin{quote}
\kn{ಅಪಘಾತದಲ್ಲಿ} \kn{ಅವನಿಗೆ} \kn{ತುಂಬಾ} \kn{ರಕ್ತಸ್ರಾವ} \kn{ಆಯಿತು}.
\end{quote}

\begin{quote}
\emph{apaghātadalli avanige tumbā raktasrāva āyitu.}
\end{quote}

\begin{quote}
He bled heavily in the accident.
\end{quote}

\kn{ಸ್ರಾವ} is a flow.

\begin{cousinweb}
\textbf{\kn{ರಕ್ತಸ್ರಾವ}} (\emph{raktasrāva}) is a learned compound built from Sanskrit \emph{rakta}, \textquotedblleft{}blood\textquotedblright{} (CDIAL 10539), and \emph{srāva}, \textquotedblleft{}a flowing\textquotedblright{} (Monier-Williams).
\end{cousinweb}

\begin{grammarlens}[title={What you've built}]
One more word for describing what you lived and felt.
\end{grammarlens}

\subsection*{Your turn}
\begin{itemize}
\raggedright
  \item \emph{Say it:} \emph{raktasrāva}
  \item \emph{Say it:} \emph{raktasrāva}, once more
  \item \emph{Say it:} \emph{apaghātadalli avanige tumbā raktasrāva āyitu.} — He bled heavily in the accident.
  \item \emph{Recall:} say the Kannada for blood pressure, then the Kannada for a heart attack, then say \emph{raktasrāva} again
\end{itemize}

\begin{tcolorbox}[breakable,colback=teal!4,colframe=teal!35!black,title={Before you move on}]
What does \kn{ರಕ್ತಸ್ರಾವ} mean? (\textquotedblleft{}Bleeding; a haemorrhage\textquotedblright{}.) Say it once more.
\end{tcolorbox}
\section[\texorpdfstring{\kn{ಪಾರ್ಶ್ವವಾಯು} (pārśvavāyu)}{ಪಾರ್ಶ್ವವಾಯು (pārśvavāyu)}]{\kn{ಪಾರ್ಶ್ವವಾಯು} (pārśvavāyu) — a stroke; paralysis}

\label{lesson:KA-C512-parsvavayu}



\begin{quote}
Before the new one: say the Kannada for a heart attack, then the Kannada for bleeding.
\end{quote}

\subsection*{You'll want to know: \kn{ಪಾರ್ಶ್ವವಾಯು}}
\textbf{\kn{ಪಾರ್ಶ್ವವಾಯು}} — \emph{pārśvavāyu} — \textquotedblleft{}a stroke; paralysis\textquotedblright{}.

\begin{quote}
\kn{ಪಾರ್ಶ್ವವಾಯುವಿನಿಂದ} \kn{ಅಜ್ಜನ} \kn{ಎಡಗೈ} \kn{ಚಲಿಸುವುದಿಲ್ಲ}.
\end{quote}

\begin{quote}
\emph{pārśvavāyuvininda ajjana eḍagai calisuvudilla.}
\end{quote}

\begin{quote}
Since his stroke, grandfather's left hand does not move.
\end{quote}

Literally side-wind: \kn{ಪಾರ್ಶ್ವ}, a side + \kn{ವಾಯು}, air, the old name for a disorder of the body.

\begin{cousinweb}
\textbf{\kn{ಪಾರ್ಶ್ವವಾಯು}} (\emph{pārśvavāyu}) is a learned compound built from Sanskrit \emph{pārśva}, \textquotedblleft{}the side, the ribs\textquotedblright{} (CDIAL 8118), and \emph{vāyu}, \textquotedblleft{}wind\textquotedblright{} (CDIAL 11544): a \textquotedblleft{}wind\textquotedblright{} that strikes one side of the body.
\end{cousinweb}

\begin{grammarlens}[title={What you've built}]
One more word for describing what you lived and felt.
\end{grammarlens}

\subsection*{Your turn}
\begin{itemize}
\raggedright
  \item \emph{Say it:} \emph{pārśvavāyu}
  \item \emph{Say it:} \emph{pārśvavāyu}, once more
  \item \emph{Say it:} \emph{pārśvavāyuvininda ajjana eḍagai calisuvudilla.} — Since his stroke, grandfather's left hand does not move.
  \item \emph{Recall:} say the Kannada for a heart attack, then the Kannada for bleeding, then say \emph{pārśvavāyu} again
\end{itemize}

\begin{tcolorbox}[breakable,colback=teal!4,colframe=teal!35!black,title={Before you move on}]
What does \kn{ಪಾರ್ಶ್ವವಾಯು} mean? (\textquotedblleft{}A stroke; paralysis\textquotedblright{}.) Say it once more.
\end{tcolorbox}
\section[\texorpdfstring{\kn{ಪ್ರಜ್ಞೆ} \kn{ತಪ್ಪು} (prajñe tappu)}{ಪ್ರಜ್ಞೆ ತಪ್ಪು (prajñe tappu)}]{\kn{ಪ್ರಜ್ಞೆ} \kn{ತಪ್ಪು} (prajñe tappu) — to faint; to pass out}

\label{lesson:KA-C512-prajne-tappu}



\begin{quote}
Before the new one: say the Kannada for bleeding, then the Kannada for a stroke.
\end{quote}

\subsection*{You'll want to know: \kn{ಪ್ರಜ್ಞೆ} \kn{ತಪ್ಪು}}
\textbf{\kn{ಪ್ರಜ್ಞೆ} \kn{ತಪ್ಪು}} — \emph{prajñe tappu} — \textquotedblleft{}to faint; to pass out\textquotedblright{}.

\begin{quote}
\kn{ಬಿಸಿಲಿನಲ್ಲಿ} \kn{ನಿಂತಿದ್ದ} \kn{ಹುಡುಗನಿಗೆ} \kn{ಪ್ರಜ್ಞೆ} \kn{ತಪ್ಪಿತು}.
\end{quote}

\begin{quote}
\emph{bisilinalli nintidda huḍuganige prajñe tappitu.}
\end{quote}

\begin{quote}
The boy who had been standing in the sun fainted.
\end{quote}

Past (it …): \kn{ಪ್ರಜ್ಞೆ} \kn{ತಪ್ಪಿತು}, \emph{prajñe tappitu}.

\kn{ಪ್ರಜ್ಞೆ} is consciousness; the one who faints takes -\kn{ಗೆ}. Coming round is \kn{ಪ್ರಜ್ಞೆ} \kn{ಬರು}.

\begin{cousinweb}
\textbf{\kn{ಪ್ರಜ್ಞೆ} \kn{ತಪ್ಪು}} (\emph{prajñe tappu}) joins Sanskrit \emph{prajñā}, \textquotedblleft{}wisdom, awareness\textquotedblright{} (CDIAL 8515), and the native verb \emph{tappu}, \textquotedblleft{}to slip, to miss\textquotedblright{}: awareness slips away.
\end{cousinweb}

\begin{grammarlens}[title={What you've built}]
One more word for describing what you lived and felt.
\end{grammarlens}

\subsection*{Your turn}
\begin{itemize}
\raggedright
  \item \emph{Say it:} \emph{prajñe tappu}
  \item \emph{Say it:} \emph{prajñe tappu}, once more
  \item \emph{Say it:} \emph{bisilinalli nintidda huḍuganige prajñe tappitu.} — The boy who had been standing in the sun fainted.
  \item \emph{Recall:} say the Kannada for bleeding, then the Kannada for a stroke, then say \emph{prajñe tappu} again
\end{itemize}

\begin{tcolorbox}[breakable,colback=teal!4,colframe=teal!35!black,title={Before you move on}]
What does \kn{ಪ್ರಜ್ಞೆ} \kn{ತಪ್ಪು} mean? (\textquotedblleft{}To faint; to pass out\textquotedblright{}.) Say it once more.
\end{tcolorbox}
